% !TeX program = pdflatex
% !BIB program = biber
%
% Plik główny — „skleja” całą pracę. Kompilacja: make (albo latexmk), wynik: build/main.pdf
%
\documentclass{pollubthesis}   % [oneside] — wydruk jednostronny (bez marginesów lustrzanych)

% Dane na stronę tytułową (i do metadanych PDF).
% Dane uzupełnione na podstawie plakatu CITS26 — sprawdź pozycje oznaczone TODO.

\autor{Karol Mac}
\nralbumu{99058}               % TODO: sprawdź (wzięte z adresu s99058@pollub.edu.pl)
\rodzajpracy{magisterska}      % „Praca dyplomowa magisterska”
\kierunek{Informatyka}
\specjalnosc{}                 % TODO: nazwa specjalności; puste {} = linia „w specjalności …” jest pomijana

% TODO: wpisz oficjalny tytuł pracy (dokładnie taki, jak w systemie / u promotora).
% Poniżej tytuł plakatu z CITS26 i jego tłumaczenie — tylko jako punkt wyjścia.
\tytul{Analiza porównawcza wybranych cech języków Java i~Rust na~podstawie programów realizujących równoważne funkcjonalności}
\tytulang{Comparative Analysis of Selected Characteristics of Java and Rust Based on Programs Implementing Equivalent Functionalities}

\promotor{dr inż. Łukasz Gałka}  % TODO: sprawdź tytuł/stopień naukowy promotora
\miejscowosc{Lublin}
\rok{\the\year}                % rok kompilacji; wpisz na sztywno, np. \rok{2027}, jeśli trzeba

\addbibresource{bibliografia.bib}

\begin{document}

% Strona tytułowa — układ z szablonu Word (Politechnika Lubelska, WEiI).
% Dane pochodzą z metadane.tex; tutaj zmieniasz tylko układ.
% Odstępy pionowe odtwarzają puste akapity z Worda (wymierzone na wydruku oryginału).
\begingroup
\thispagestyle{empty}
\setstretch{1}
\setlength{\parindent}{0pt}
\hyphenpenalty=10000
\exhyphenpenalty=10000
\sloppy
\hbadness=10000  % wyjustowany tytuł bez dzielenia wyrazów — rozciągnięte odstępy są zamierzone (jak w Wordzie)
\lineskiplimit=-\maxdimen  % odstępy zawsze liczone od linii bazowych (także po dużym napisie)
\fontsize{12}{13.8}\selectfont

\includegraphics[width=8.75cm]{logo-weii}\par
\vspace{27pt}
{\raggedright\sffamily\fontsize{46}{52.9}\selectfont\color{pollubtekst}%
  Praca dyplomowa \PracaRodzaj\par}
\vspace{18.8pt}
{\color{pollubtekst}%
  Na kierunku \PracaKierunek\par
  \ifdefempty{\PracaSpecjalnosc}{}{w specjalności \PracaSpecjalnosc\par}}
\vspace{36.3pt}
{\fontsize{15}{17.25}\selectfont\bfseries \PracaTytul\par}
\vspace{17.25pt}
{\fontsize{15}{17.25}\selectfont\bfseries \foreignlanguage{english}{\PracaTytulAng}\par}
\vspace{52.7pt}
{\fontsize{16}{18.4}\selectfont\bfseries \PracaAutor\par}
\vspace{12.4pt}
{\color{pollubtekst}Nr albumu: \PracaNrAlbumu\par}
\vspace{21.9pt}
{\color{pollubtekst}Promotor: \PracaPromotor\par}
\vfill
{\fontsize{10}{11.5}\selectfont\hspace*{0.245cm}\PracaMiejscowosc{} \PracaRok\par}
\vspace{4.5pt}
\endgroup
\cleardoublepage  % przy druku dwustronnym: pusta strona 2 (rewers strony tytułowej)

\spistresci
\begin{streszczenie}
% TODO: treść streszczenia po polsku.

\slowakluczowe{% TODO: słowa kluczowe, po przecinku
}
\end{streszczenie}

\begin{otherlanguage}{english}  % angielskie dzielenie wyrazów
\begin{streszczenie}[Abstract]
% TODO: abstract in English.

\slowakluczowe[Key words]{% TODO: keywords, comma separated
}
\end{streszczenie}
\end{otherlanguage}


\chapter{Wstęp}
\label{chap:wstep}

% TODO: wprowadzenie do tematu, motywacja, krótki opis układu pracy.

\chapter{Cel i zakres pracy}
\label{chap:cel-i-zakres}

\section{Cel pracy}
\label{sec:cel-pracy}

% TODO: cel główny i cele szczegółowe.

\section{Zakres pracy}
\label{sec:zakres-pracy}

% TODO

\chapter{Podstawy teoretyczne}
\label{chap:podstawy}

% Szkielet — zmień tytuły i liczbę podrozdziałów według potrzeb.

\section{Tytuł podrozdziału}

\subsection{Tytuł punktu}

% TODO

\subsection{Tytuł punktu}

% TODO

\section{Tytuł podrozdziału}

\subsection{Tytuł punktu}

% TODO

\chapter{Metodologia badań}
\label{chap:metodologia}

\section{Środowisko badawcze i dane testowe}

% TODO

\section{Implementacja pierwszego rozwiązania}

% TODO

\section{Implementacja drugiego rozwiązania}

% TODO

\section{Metody oceny i metryki}

% TODO

\section{Procedura eksperymentu}

% TODO

\chapter{Wyniki}
\label{chap:wyniki}

\section{Tytuł podrozdziału}

% TODO

\section{Tytuł podrozdziału}

% TODO

\chapter{Wnioski}
\label{chap:wnioski}

% TODO


% ŚCIĄGA — przykłady użycia szablonu. Usuń % ŚCIĄGA — przykłady użycia szablonu. Usuń % ŚCIĄGA — przykłady użycia szablonu. Usuń \include{rozdzialy/99-przyklady} z main.tex przed oddaniem.
\chapter{Przykłady użycia szablonu (do usunięcia)}
\label{chap:przyklady}

\section{Tekst, cytowania i odwołania}

Zwykły akapit jest wyjustowany, ma wcięcie 1,27~cm i~interlinię 1,5. Cytowanie pojedynczej
pozycji wygląda tak~\cite{przyklad-ksiazka}, kilku naraz tak~\cite{przyklad-artykul,
przyklad-konferencja,przyklad-www}. Kolejność w~bibliografii wynika z~kolejności pierwszego
cytowania. Odwołania do innych miejsc: rozdział~\ref{chap:wstep}, rysunek~\ref{fig:przyklad},
tabela~\ref{tab:przyklad}, wzór~\eqref{eq:przyklad}, listing~\ref{lst:java}\footnote{Przypis
dolny wstawia się poleceniem \texttt{\textbackslash footnote}.}.

Jednoliterowe spójniki i~przyimki (a, i, o, u, w, z) wiąże się z~następnym słowem tyldą:
\texttt{w\textasciitilde{}pracy}, aby nie zostawały na końcu wiersza. Liczby z~przecinkiem
dziesiętnym: \num{1234.5} albo z~jednostką: \SI{12.5}{\milli\second}.

\section{Listy}

\begin{itemize}
  \item pierwszy element listy punktowanej,
  \item drugi element:
  \begin{itemize}
    \item element drugiego poziomu,
    \item kolejny element drugiego poziomu,
  \end{itemize}
  \item trzeci element.
\end{itemize}

\begin{enumerate}
  \item pierwszy krok,
  \item drugi krok.
\end{enumerate}

\section{Rysunek, wykres i tabela}

Rysunek wstawiany jest z~katalogu \texttt{grafika/} (najlepiej w~PDF lub PNG). Podpis jest pod
rysunkiem, numeracja jest ciągła w~całej pracy.

\begin{figure}[htbp]
  \centering
  \includegraphics[width=0.5\textwidth]{example-image}
  \caption{Przykładowy rysunek}
  \label{fig:przyklad}
\end{figure}

Wykresy mają osobną numerację (w~oryginale „Wykres 1”, „Wykres 2”, …). Jeśli wolisz traktować
je jak rysunki, użyj zwykłego \texttt{figure}.

\begin{wykres}[htbp]
  \centering
  \includegraphics[width=0.5\textwidth]{example-image-a}
  \caption{Przykładowy wykres}
  \label{wyk:przyklad}
\end{wykres}

Tabela w~stylu oryginału: pełna siatka, szary wiersz nagłówka, wyśrodkowane komórki, podpis pod
tabelą (aby przenieść podpis nad tabelę, przesuń \texttt{\textbackslash caption} przed
\texttt{tabular}).

\begin{table}[htbp]
  \centering
  \begin{tabular}{|c|c|c|}
    \hline
    \rowcolor{tabnaglowek}
    \textbf{Kolumna 1} & \textbf{Kolumna 2} & \textbf{Kolumna 3} \\
    \hline
    A & \num{0.95} & \num{120} \\
    \hline
    B & \num{0.87} & \num{340} \\
    \hline
  \end{tabular}
  \caption{Przykładowa tabela}
  \label{tab:przyklad}
\end{table}

\section{Wzory i listingi}

\begin{equation}
  F_1 = 2 \cdot \frac{\mathit{precision} \cdot \mathit{recall}}{\mathit{precision} + \mathit{recall}}
  \label{eq:przyklad}
\end{equation}

\begin{lstlisting}[language=Java, caption={Przykładowy kod w~języku Java}, label={lst:java}]
public class Main {
    public static void main(String[] args) {
        // komentarz z polskimi znakami: zażółć gęślą jaźń
        System.out.println("Witaj!");
    }
}
\end{lstlisting}

\begin{lstlisting}[language=Rust, caption={Przykładowy kod w~języku Rust}, label={lst:rust}]
fn main() {
    let liczby: Vec<u64> = (1..=10).collect();
    let suma: u64 = liczby.iter().sum();
    println!("Suma: {}", suma);
}
\end{lstlisting}

Kod można też wczytać wprost z~pliku źródłowego:

\begin{verbatim}
\lstinputlisting[language=Rust, caption={Opis}]{../rust/src/main.rs}
\end{verbatim}
 z main.tex przed oddaniem.
\chapter{Przykłady użycia szablonu (do usunięcia)}
\label{chap:przyklady}

\section{Tekst, cytowania i odwołania}

Zwykły akapit jest wyjustowany, ma wcięcie 1,27~cm i~interlinię 1,5. Cytowanie pojedynczej
pozycji wygląda tak~\cite{przyklad-ksiazka}, kilku naraz tak~\cite{przyklad-artykul,
przyklad-konferencja,przyklad-www}. Kolejność w~bibliografii wynika z~kolejności pierwszego
cytowania. Odwołania do innych miejsc: rozdział~\ref{chap:wstep}, rysunek~\ref{fig:przyklad},
tabela~\ref{tab:przyklad}, wzór~\eqref{eq:przyklad}, listing~\ref{lst:java}\footnote{Przypis
dolny wstawia się poleceniem \texttt{\textbackslash footnote}.}.

Jednoliterowe spójniki i~przyimki (a, i, o, u, w, z) wiąże się z~następnym słowem tyldą:
\texttt{w\textasciitilde{}pracy}, aby nie zostawały na końcu wiersza. Liczby z~przecinkiem
dziesiętnym: \num{1234.5} albo z~jednostką: \SI{12.5}{\milli\second}.

\section{Listy}

\begin{itemize}
  \item pierwszy element listy punktowanej,
  \item drugi element:
  \begin{itemize}
    \item element drugiego poziomu,
    \item kolejny element drugiego poziomu,
  \end{itemize}
  \item trzeci element.
\end{itemize}

\begin{enumerate}
  \item pierwszy krok,
  \item drugi krok.
\end{enumerate}

\section{Rysunek, wykres i tabela}

Rysunek wstawiany jest z~katalogu \texttt{grafika/} (najlepiej w~PDF lub PNG). Podpis jest pod
rysunkiem, numeracja jest ciągła w~całej pracy.

\begin{figure}[htbp]
  \centering
  \includegraphics[width=0.5\textwidth]{example-image}
  \caption{Przykładowy rysunek}
  \label{fig:przyklad}
\end{figure}

Wykresy mają osobną numerację (w~oryginale „Wykres 1”, „Wykres 2”, …). Jeśli wolisz traktować
je jak rysunki, użyj zwykłego \texttt{figure}.

\begin{wykres}[htbp]
  \centering
  \includegraphics[width=0.5\textwidth]{example-image-a}
  \caption{Przykładowy wykres}
  \label{wyk:przyklad}
\end{wykres}

Tabela w~stylu oryginału: pełna siatka, szary wiersz nagłówka, wyśrodkowane komórki, podpis pod
tabelą (aby przenieść podpis nad tabelę, przesuń \texttt{\textbackslash caption} przed
\texttt{tabular}).

\begin{table}[htbp]
  \centering
  \begin{tabular}{|c|c|c|}
    \hline
    \rowcolor{tabnaglowek}
    \textbf{Kolumna 1} & \textbf{Kolumna 2} & \textbf{Kolumna 3} \\
    \hline
    A & \num{0.95} & \num{120} \\
    \hline
    B & \num{0.87} & \num{340} \\
    \hline
  \end{tabular}
  \caption{Przykładowa tabela}
  \label{tab:przyklad}
\end{table}

\section{Wzory i listingi}

\begin{equation}
  F_1 = 2 \cdot \frac{\mathit{precision} \cdot \mathit{recall}}{\mathit{precision} + \mathit{recall}}
  \label{eq:przyklad}
\end{equation}

\begin{lstlisting}[language=Java, caption={Przykładowy kod w~języku Java}, label={lst:java}]
public class Main {
    public static void main(String[] args) {
        // komentarz z polskimi znakami: zażółć gęślą jaźń
        System.out.println("Witaj!");
    }
}
\end{lstlisting}

\begin{lstlisting}[language=Rust, caption={Przykładowy kod w~języku Rust}, label={lst:rust}]
fn main() {
    let liczby: Vec<u64> = (1..=10).collect();
    let suma: u64 = liczby.iter().sum();
    println!("Suma: {}", suma);
}
\end{lstlisting}

Kod można też wczytać wprost z~pliku źródłowego:

\begin{verbatim}
\lstinputlisting[language=Rust, caption={Opis}]{../rust/src/main.rs}
\end{verbatim}
 z main.tex przed oddaniem.
\chapter{Przykłady użycia szablonu (do usunięcia)}
\label{chap:przyklady}

\section{Tekst, cytowania i odwołania}

Zwykły akapit jest wyjustowany, ma wcięcie 1,27~cm i~interlinię 1,5. Cytowanie pojedynczej
pozycji wygląda tak~\cite{przyklad-ksiazka}, kilku naraz tak~\cite{przyklad-artykul,
przyklad-konferencja,przyklad-www}. Kolejność w~bibliografii wynika z~kolejności pierwszego
cytowania. Odwołania do innych miejsc: rozdział~\ref{chap:wstep}, rysunek~\ref{fig:przyklad},
tabela~\ref{tab:przyklad}, wzór~\eqref{eq:przyklad}, listing~\ref{lst:java}\footnote{Przypis
dolny wstawia się poleceniem \texttt{\textbackslash footnote}.}.

Jednoliterowe spójniki i~przyimki (a, i, o, u, w, z) wiąże się z~następnym słowem tyldą:
\texttt{w\textasciitilde{}pracy}, aby nie zostawały na końcu wiersza. Liczby z~przecinkiem
dziesiętnym: \num{1234.5} albo z~jednostką: \SI{12.5}{\milli\second}.

\section{Listy}

\begin{itemize}
  \item pierwszy element listy punktowanej,
  \item drugi element:
  \begin{itemize}
    \item element drugiego poziomu,
    \item kolejny element drugiego poziomu,
  \end{itemize}
  \item trzeci element.
\end{itemize}

\begin{enumerate}
  \item pierwszy krok,
  \item drugi krok.
\end{enumerate}

\section{Rysunek, wykres i tabela}

Rysunek wstawiany jest z~katalogu \texttt{grafika/} (najlepiej w~PDF lub PNG). Podpis jest pod
rysunkiem, numeracja jest ciągła w~całej pracy.

\begin{figure}[htbp]
  \centering
  \includegraphics[width=0.5\textwidth]{example-image}
  \caption{Przykładowy rysunek}
  \label{fig:przyklad}
\end{figure}

Wykresy mają osobną numerację (w~oryginale „Wykres 1”, „Wykres 2”, …). Jeśli wolisz traktować
je jak rysunki, użyj zwykłego \texttt{figure}.

\begin{wykres}[htbp]
  \centering
  \includegraphics[width=0.5\textwidth]{example-image-a}
  \caption{Przykładowy wykres}
  \label{wyk:przyklad}
\end{wykres}

Tabela w~stylu oryginału: pełna siatka, szary wiersz nagłówka, wyśrodkowane komórki, podpis pod
tabelą (aby przenieść podpis nad tabelę, przesuń \texttt{\textbackslash caption} przed
\texttt{tabular}).

\begin{table}[htbp]
  \centering
  \begin{tabular}{|c|c|c|}
    \hline
    \rowcolor{tabnaglowek}
    \textbf{Kolumna 1} & \textbf{Kolumna 2} & \textbf{Kolumna 3} \\
    \hline
    A & \num{0.95} & \num{120} \\
    \hline
    B & \num{0.87} & \num{340} \\
    \hline
  \end{tabular}
  \caption{Przykładowa tabela}
  \label{tab:przyklad}
\end{table}

\section{Wzory i listingi}

\begin{equation}
  F_1 = 2 \cdot \frac{\mathit{precision} \cdot \mathit{recall}}{\mathit{precision} + \mathit{recall}}
  \label{eq:przyklad}
\end{equation}

\begin{lstlisting}[language=Java, caption={Przykładowy kod w~języku Java}, label={lst:java}]
public class Main {
    public static void main(String[] args) {
        // komentarz z polskimi znakami: zażółć gęślą jaźń
        System.out.println("Witaj!");
    }
}
\end{lstlisting}

\begin{lstlisting}[language=Rust, caption={Przykładowy kod w~języku Rust}, label={lst:rust}]
fn main() {
    let liczby: Vec<u64> = (1..=10).collect();
    let suma: u64 = liczby.iter().sum();
    println!("Suma: {}", suma);
}
\end{lstlisting}

Kod można też wczytać wprost z~pliku źródłowego:

\begin{verbatim}
\lstinputlisting[language=Rust, caption={Opis}]{../rust/src/main.rs}
\end{verbatim}
    % ŚCIĄGA z przykładami (rysunek, tabela, listing, cytowanie) — usuń przed oddaniem

\printbibliography[heading=bibintoc]

% Spisy — w oryginale ich nie było; odkomentuj, jeśli promotor ich wymaga:
% \listoffigures
% \listoftables
% \listof{wykres}{Spis wykresów}
% \lstlistoflistings

\end{document}
